\section{Benchmark Setup}\label{sec-benchmark}
We study binary post-wildfire debris-flow (PFDF) occurrence prediction at the storm-watershed-event level~\cite{staley2016}. Each sample corresponds to a paired observation of a burned watershed segment and a rainfall storm event ${(\mathbf{x}_i, y_i)}$. For each post-fire storm affecting a monitored burned watershed segment, the dataset records storm forcing conditions, watershed and burn attributes, and a binary label $y_i \in \{0,1\}$ indicating whether a debris flow was observed during the storm response. The input feature $\mathbf{x}_i$ consists of meteorological features
(e.g., rainfall intensity and storm duration), wildfire features (e.g., burn severity and dNBR), and geologic/geomorphic
features (e.g., terrain slope and soil erodibility).

\subsection{Design Space} 

\subsubsection{Dataset.} We leverage the post-fire debris-flow database released by~\cite{staley2016}, where each entry represents a monitored storm response for a burned watershed segment, including both debris-flow and non-debris-flow outcomes. The dataset contains 1,243 storm-watershed events from 33 fires across 10 geographic zones in the western United States, with 316 positive cases (25.4\%).
%
For consistent evaluation, we use a complete-case subset by removing events with missing sub-hourly rainfall variables required by baseline models, yielding 1,131 events with 292 positives (25.8\%). Each event includes storm summaries, burn severity, soil and terrain attributes, and spatial-temporal indicators.

\subsubsection{Feature Setup.}
The original feature table follows the USGS post-fire debris-flow database~\cite{staley2016}. Each storm-watershed observation contains a binary debris-flow response with a total of 26 attributes spanning meteorology, wildfire, and geology~\cite{staley2016updated, gartner2014empirical}. 
% Storm variables characterize rainfall forcing through duration, accumulation, average intensity, and short-window peak rainfall statistics at 15, 30, and 60 minutes. Burn and watershed attributes include contributing area, dNBR-based burn severity, steep high-severity burn proportion, soil erodibility (KF factor), geographic zone, storm month, and rain-gauge distance. 
These variables form the core feature for the USGS-deployed model.

We further augment the baseline event attributes by linking each record to additional physical and meteorological context. Terrain attributes (e.g., mean slope and aspect), are derived from USGS 3DEP elevation products based on the geo-coordinates of each event\footnote{\url{https://www.usgs.gov/3d-elevation-program}}. Antecedent hydrologic conditions are represented by NLDAS-2 soil-moisture summaries and time since fire. In addition, we match storm events to National Weather Service forecast products\footnote{\url{https://mesonet.agron.iastate.edu/api/}}, including Area Forecast Discussions and Hazardous Weather Outlooks. We then prompt an LLM (Table~\ref{tab:prompt_weather_example}) to convert these narrative forecasts into structured PFDF weather indicators, such as precipitation threat, extreme rainfall potential, and overall PFDF risk.
% Overall, the augmented table extends the original USGS features with terrain, hydrologic recovery, and forecast-derived meteorological information.


\subsubsection{Baseline and Model Families.}\label{sec:baselines}
Following prior studies~\cite{staley2016, kern2017ml_pfdf, nikolopoulos2018pfdf_rf}, we benchmark a fixed operational reference model and a suite of trainable learning baselines. Physics logistic regression (P-LR) denotes the published USGS operational model with frozen coefficients~\cite{staley2016}. The trainable baselines include logistic regression (LR), classification and regression tree (CART)~\cite{kern2017ml_pfdf}, support vector machine (SVM)~\cite{kern2017ml_pfdf}, neural network (NN)~\cite{kern2017ml_pfdf}, random forest (RF)~\cite{nikolopoulos2018pfdf_rf}, and histogram gradient boosting (HGB). For each baseline, we also evaluate an augmented-feature variant with terrain attributes, antecedent hydrologic conditions, and LLM-processed PFDF weather indicators, denoted by $\mathrm{LR}^{*}$, $\mathrm{CART}^{*}$, $\mathrm{SVM}^{*}$, $\mathrm{NN}^{*}$, $\mathrm{RF}^{*}$, and $\mathrm{HGB}^{*}$.

\subsubsection{Evaluation Metric.} Our benchmark computes a comprehensive set of binary-classification metrics,
including accuracy, precision, recall, false-alarm behavior, and Threat Score~\cite{staley2016, kern2017ml_pfdf, nikolopoulos2018pfdf_rf}. Following PFDF evaluation used for USGS deployment model~\cite{staley2016}, we report Threat Score (TS): $\mathrm{TS}=\frac{\mathrm{TP}}{\mathrm{TP}+\mathrm{FP}+\mathrm{FN}}$. TS is well-suited to imbalanced hazard-detection problems because it directly
rewards correct debris-flow detections while penalizing both missed events (false negatives) and false alarms (false positives), without allowing the large number of true negatives to dominate the evaluation.


\begin{table}[t!]
\centering
\caption{The LLM is constrained to fill a predefined JSON schema with 18 numeric PFDF-relevant weather indicators and one concise textual summary. These fields decompose broad forecast signals, including rainfall intensity, storm duration or repetition, flash-flood language, convection, terrain enhancement, antecedent wetness, burn-scar sensitivity, and overall PFDF weather risk.}
\label{tab:prompt_weather_example}
\scriptsize
\begin{tcolorbox}[
    colback=white,
    colframe=green!40!black,
    boxrule=1pt,
    title={LLM-based Processing to Structure Area Forecast Discussion Text Corpus},
    colbacktitle=green!40!black,
    coltitle=white
]

\textbf{Input AFD:} \texttt{AFDLOX@2009-12-09T22:01}

\textbf{Forecast excerpt:}
``Relatively cold temperatures will continue tonight, with slight warming expected
through the week. A storm system will bring a chance of light to moderate rain on
Thursday into early Friday. Another storm system is expected to affect the area
over the weekend, with the potential for more significant rainfall. This storm
produces significant orographic rainfall, with preliminary totals of 0.5 to 1.5
inches for coastal and valley areas and 1.5 to 2.5 inches in the foothills and
mountains.''

\textbf{Extraction instruction:}
Read the AFD as an operational meteorologist and post-wildfire debris-flow hazard analyst. Convert only the PFDF-relevant weather signal into a compact and structured record. Use only the information stated in the forecast text.

\textbf{Structured output:}

\emph{Summary:} ``Two storm systems are expected with moderate to significant
rainfall, especially in foothills and mountains due to orographic enhancement.''

\begin{center}
\scriptsize
\setlength{\tabcolsep}{6pt}
\renewcommand{\arraystretch}{1.05}
\resizebox{\linewidth}{!}{%
\begin{tabular}{@{}>{\ttfamily}l>{\ttfamily}l@{}}
precipitation\_signal=1 & localized\_extreme\_signal=0 \\
short\_duration\_intensity\_signal=0 & mentions\_heavy\_rain=0 \\
storm\_duration\_signal=1 & mentions\_thunderstorm=0 \\
training\_repetition\_signal=0 & mentions\_flash\_flood=0 \\
antecedent\_wetness\_signal=0 & mentions\_burn\_scar\_or\_debris\_flow=0 \\
convective\_signal=0 & mentions\_training\_or\_repeating\_cells=0 \\
orographic\_enhancement\_signal=1 & mentions\_slow\_storm\_motion=0 \\
explicit\_flash\_flood\_signal=0 & mentions\_monsoon\_or\_tropical\_moisture=0 \\
burn\_scar\_sensitivity\_signal=0 & overall\_pfdf\_weather\_risk=1 \\
\end{tabular}%
}
\end{center}
\end{tcolorbox}
\end{table}

\subsubsection{Hyperparameter.}
We evaluate four zone stratified train/\allowbreak validation/\allowbreak test ratios:
50/25/25, 60/20/20, 70/15/15, and 80/10/10, using five split seeds. For each ratio and split seed, hyperparameters are grid searched by validation Threat Score over learning rate $\{10^{-4},\allowbreak 3{\times}10^{-4},\allowbreak 10^{-3},\allowbreak 3{\times}10^{-3}\}$,
dropout $\{0.0,0.1,0.2,0.3,0.5\}$, hidden width
$\{32,64,128,256,512\}$, and weight decay $\{0,10^{-4}\}$.
% For neural models, the validation Threat Score is averaged over initialization seeds 42-46 for hyperparameter selection, and the selected configuration is evaluated over initialization seeds 42-51. 
All neural models are trained with Adam, binary cross-entropy, and early stopping with patience 50 on validation Threat Score.

% 
\subsection{RL-based Feature Selection}\label{sec-RL}
Given trained PFDF prediction models $\mathcal{F}$, we learn a binary feature-removal mask $\mathbf{M}\in\{0,1\}^{d}$ to identify features responsible for separating positive and negative PFDF events. Here, $\mathbf{M}_{j}=1$ suppresses feature $j$, while $\mathbf{M}_{j}=0$ retains it. Masked features are replaced by reference values: $\tilde{\mathbf{x}}_{\mathbf{M}}
= (\mathbf{1}-\mathbf{M})\odot \mathbf{x}
+ \mathbf{M}\odot \mathbf{b}$ where $\mathbf{b}$ is a baseline feature vector estimated from the training set. Standardized continuous features use $0$ (training-set mean after normalization), while the ordinal zone feature uses the most frequent training-zone value. For a model $f\in\mathcal{F}$, we measure the predictive separation between positive and negative samples under mask $\mathbf{M}$ using the standardized mean gap:
$D_f(\mathbf{M})=
\frac{
\left|
\mathbb{E}_{\mathbf{x}^{+}\sim \mathcal{D}^{+}}
f(\tilde{\mathbf{x}}^{+}_{\mathbf{M}})
-
\mathbb{E}_{\mathbf{x}^{-}\sim \mathcal{D}^{-}}
f(\tilde{\mathbf{x}}^{-}_{\mathbf{M}})
\right|
}{
\sqrt{
\mathrm{Var}_{\mathbf{x}^{+}\sim \mathcal{D}^{+}}
f(\tilde{\mathbf{x}}^{+}_{\mathbf{M}})
+
\mathrm{Var}_{\mathbf{x}^{-}\sim \mathcal{D}^{-}}
f(\tilde{\mathbf{x}}^{-}_{\mathbf{M}})
+\epsilon
}}$.
We then optimize masks that maximally collapse this separation relative to the unmasked input: $\mathbf{M}^{*}
= \arg\max_{\mathbf{M}}\mathbb{E}_{f\sim\mathcal{F}}
[1-\frac{D_f(\mathbf{M})}{D_f(\boldsymbol{0})}
] -\lambda \|\mathbf{M}\|_{0}$. This objective identifies features whose removal makes positive and negative predictions less distinguishable, implying their importance for PFDF discrimination, while the sparsity term avoids suppressing excessive features. We optimize the discrete masking objective using a REINFORCE-style policy gradient~\cite{liu2021automated, fan2020autofs, fan2021interactive, sutton1999policy}.




% \begin{table}[t!]
% \centering
% \caption{Dataset used for the controlled benchmark.}
% \label{tab-data}
% \begin{tabular}{lr}
% \toprule
% \textbf{Attribute} & \textbf{Value} \\
% \midrule
% Raw monitored storm--watershed events & 1,243 \\
% Positive responses, raw & 316 (25.4\%) \\
% Fires & 33 \\
% Geographic zones & 10 \\
% Complete-case benchmark events & 1,131 \\
% Positive responses, complete case & 292 (25.8\%) \\
% \bottomrule
% \end{tabular}
% \end{table}

% \begin{table}[t]
% \centering
% \caption{Benchmark model families and the comparison each one supports.}
% \label{tab:model_families}
% \small
% \begin{tabularx}{\linewidth}{p{1.8cm}X}
% \toprule
% \textbf{Model} & \textbf{Role in the benchmark} \\
% \midrule
% P-LR & Fixed USGS logistic-regression reference with frozen published coefficients. \\
% NN & Structured neural baseline using only tabular storm, watershed, burn, season, and zone inputs. \\
% NN+W & Tests whether dense pre-event forecast-text embeddings add predictive value. \\
% NN+SW & Tests whether compact structured weather indicators preserve forecast signal more efficiently. \\
% NN+AUG & Tests whether terrain, antecedent soil moisture, and recovery-time covariates improve the tabular baseline. \\
% NN+selected SW+AUG & Tests whether the strongest weather and physical covariates can be fused in a simple MLP. \\
% NN+G & Tests whether a physics-informed graph constraint improves over an unconstrained augmented MLP. \\
% \bottomrule
% \end{tabularx}
% \end{table}


